\section*{政策评价问题与观察边界}
公共支持对充电基础设施的评价不能只问交付了多少设施。建设交付、覆盖扩展、持续服务和活动增长属于不同的评价对象，也可能在不同时间和空间尺度上变化。设施可以先于活动进入低吞吐地区，承担覆盖和可达性任务；同样的活动强度下降，也可能来自活动继续增长但落后于设施扩张，或来自活动绝对收缩。把这些情形压缩为“活动除以设施”的单一比值，会掩盖建设目标与运营表现之间的差别。本文据此考察活动是否跟上设施、这种差距如何分布并持续、哪些城市构成网络变化，以及地方政策把哪些建设后的责任落实为可观察的运营要求。

分析使用2016—2023年的城市—周充电记录，并连接当地适用政策文本。记录来自一个运营商网络的脱敏公开档案\citep{ma2026}，因此城市和设施范围描述的是该观测网络，而不是全国公共充电基础设施。334个城市构成2,672个城市—年份的登记框架；其中293个城市有设施记录，123个城市至少一年有正活动。可观察城市从18个增加到293个，说明档案覆盖范围扩展，不能解释为全国设施从这些城市开始建设。地图中有333个可靠地理匹配城市；一个无法可靠分配边界的组合名称仍作为统计记录保留。年度模型、季度模型、政策内容和共同窗口使用不同但相互重叠的样本，分别回答规模、时间和责任问题，并非一个逐层淘汰的单一主样本。来源、匹配和样本角色见补充材料第~\siref{sec:si-sources}~节。

这种范围界定也决定了结果的解释方式：活动总量与设施总量分别由各城市的相应指标相加得到，总体强度则是二者之比；等城市统计回答城市之间的典型变化，固定成员窗口回答同一批城市是否改变。三者不能互相充当全国估计，也不能用一个层级的增长率推断另一个层级的覆盖或需求。政策文本则提供责任的制度语境，而不是把条款出现次数直接当作实施强度。

\section*{测量：设施、活动与强度}
兼容周要求源日期有效且设施指标为正。$Q$是兼容周的活动总和，$E$是同一组周的累计设施指标之和，$W$是兼容周数；由此定义每兼容周活动$A=Q/W$、平均设施$N=E/W$以及设施相对活动强度$I=A/N=Q/E$。兼容周内的编码零被保留；没有兼容周则保持缺失。编码零可能涉及实际零活动、报送或匿名化处理，不能一概解释为停运。设施指标也不等于实时可用容量，所以本文使用“记录活动”和“观察设施”，不把它们改写为完整城市需求或工程利用率。变量定义和支持范围见补充材料第~\siref{sec:si-sources}~节，总体分布见主文表~\mainref{tab:sources}。

这些量都按各自的年度或季度期间定义。$W$既是覆盖条件，也是活动暴露量；因此不应把$Q$直接理解为在完整日历上的需求总量。$E$是兼容周内累计设施指标形成的暴露量，$N$则表示平均观测设施规模，不能进一步推导额定功率、在线容量或新增建设日期。零值和缺失的区别贯穿固定窗口、分解和模型样本：零值可以进入相应统计，缺失则不被补成零。

\section*{设施扩张与活动增长}
固定36城的2019—2023年窗口给出最直接的比较：活动合计增长\AChange\%，设施增长\NChange\%，设施加权强度下降\IDecline\%。36城的设施均增加，但活动轨迹分为三类：\LongFast{}城的活动增长快于设施，\LongSlow{}城的活动与设施都增长但活动跟进较慢，\LongContract{}城则出现设施扩张而活动收缩。因而长期强度下降主要表示扩张速度不同，不能概括为所有城市的活动都减少。

分类以活动或设施增长为零、以及两者增长相等为边界，并按预设的小数容差处理精确持平。八个正活动共同窗口都显示设施合计增长快于活动；四个相邻的正活动增长窗口（2019—2020、2020—2021、2021—2022和2022—2023）固定成员数分别为39、55、77和93。在最近的2022—2023年，93个持续正活动城市的活动增长\RecentA\%，设施增长\RecentN\%，总体强度下降\RecentIDecline\%。这组窗口要求纳入城市在所需年份均有正活动和设施记录，回答的是同一类可持续正记录城市中活动与设施的相对增长问题。窗口的完整数值见主文表~\mainref{tab:growth_windows}，固定路径和分布见主文图~\mainref{fig:growth}。

允许编码零的端点比较（主表~\mainref{tab:recorded_states}）单独回答记录状态和较广观察覆盖的问题。四个相邻窗口的共同成员数分别为47、70、167和225；在最近的2022—2023年，93个两端均为正活动的$++$城市，117个两端均为编码零的$00$城市，另有10个$0+$和5个$+0$转换。117个$00$城市占城市数52\%，但只占期末设施份额3.47\%；93个$++$城市占期末设施的95.04\%。225城的活动和设施分别增长15.32\%和45.70\%，总体强度下降20.85\%。活动变化分解中，$++$城市贡献+14.907个基期指数点，$0+$贡献+0.418，$+0$贡献$-0.003$，$00$贡献为零。这个分解说明城市数量、设施规模和活动总量并不等价，也不把编码状态转化为新设施投运、用户进入或停运的证明。完整端点和状态诊断见补充材料第~\siref{sec:si-descriptive}~节。

持续性需要同城而非单期分类来判断。52个可分类城市中19城保持同类；早期活动增长快于设施的27城，后期10城仍属较快增长，10城转为较慢正增长，7城转为活动收缩，即17/27发生了增长类型变化。单个时窗的较快增长标签因而不能代替持续表现；评价应同时跟踪活动与设施的相对增速，以及这种关系是否在后续时段保持。$Other$是非有序的边界或剩余状态，不表示第四个优劣等级。转移图见主文图~\mainref{fig:transitions}，52城与73城的完整矩阵见补充材料第~\siref{sec:si-descriptive}~节。

状态分解的贡献以该窗口全体成员的基期活动为分母，所以贡献指数点可以相加，却不是各类别占净增长的百分比。它所保留的是“哪些城市的记录状态变化构成了总活动变化”，不是对用户进入、退出或设施启用的因果归因。把这条边界保留下来，才能同时解释持续正值城市承担的规模和编码零城市所代表的覆盖不确定性。

\section*{覆盖广度、设施规模与总体强度}
城市数、设施份额和网络总体强度描述的是不同问题。总体强度为城市活动合计与设施合计之比，也等于以设施份额为权重的城市强度均值；活动合计本身不再额外按设施加权。2020—2021年，55城中36城的城内强度上升，但网络总体强度仍下降约\MiddleIDecline\%。设施较多的城市在总体强度中权重更大；同时，设施份额变化可能抵消或放大城内强度变化。因此，一个城市强度上升不必然对总体强度变化作出正贡献，强度下降也不必然对应负的加权净贡献。

总体强度变化可按城市内部强度项和设施权重项分解。2020—2021年两项分别为\MiddleWithin{}和\MiddleComposition{}，2022—2023年为\RecentWithin{}和\RecentComposition{}；它们相加得到相应窗口的总体强度变化。固定交集52城并重新归一化端点权重后，负净贡献城市由\CommonEarlyNegativeN{}增至\CommonRecentNegativeN{}，前五项占负贡献绝对总量的比例由\CommonEarlyTopFive\%降至\CommonRecentTopFive\%，HHI由\CommonEarlyHHI{}降至\CommonRecentHHI{}。城市贡献和网络分解见主文表~\mainref{tab:contributions}和主文图~\mainref{fig:contributions}；完整排名与集中度诊断见补充材料第~\siref{sec:si-descriptive}~节。这是对总体强度变化组成和贡献集中度的会计描述，不是对需求、设备配置或政策行为机制的识别。

因此，225城的观察覆盖与93城的设施规模集中可以同时成立：前者说明观察覆盖更广，后者说明观察设施规模主要集中在持续正记录城市。若只报告城市数量，会低估规模差异；若只报告设施加权总体强度，又会看不见低设施份额城市的状态变化。固定交集比较只能说明这些特定窗口内的分布重组，不能外推为全国城市的长期趋势。

\section*{持续性：正记录周与强度}
设施相对强度还可分解为正记录频率、正记录周的活动规模和设施规模。令$p$为兼容周中正活动记录的比例，$q^{+}$为正记录周的平均活动，则
\[
A=pq^{+},\qquad I=\frac{pq^{+}}{N},\qquad
\Delta\ln I=\Delta\ln p+\Delta\ln q^{+}-\Delta\ln N.
\]
等城市平均的分量可以相加，中位数只描述各分布位置，不能与其他中位数相加。

在长期36城中，等城市平均对数强度变化为\LongMeanI{}，正记录比例、正周活动和设施分量分别为+\LongMeanP{}、+\LongMeanQ{}和$-\LongMeanN{}$；编码零周总数反而从\LongZeroStart{}降至\LongZeroEnd{}。近期窗口的等城市平均对数强度为+\RecentMeanI{}，但中位数为负，设施加权总体强度也下降。低基期和记录频率的变化可以造成很大的比例差异：眉山的正记录周从\MeishanPosStart{}增至\MeishanPosEnd{}，正周均值从\MeishanQStart{}升至\MeishanQEnd{}；临沂则从\LinyiPosStart{}个正周、\LinyiQStart{}的均值降至\LinyiPosEnd{}个正周、\LinyiQEnd{}。这些记录基础说明，等权平均、中位数和设施加权总体强度的方向差异不能由一个统计量互相替代；记录分解见主文图~\mainref{fig:recording}，完整记录频率和端点诊断见补充材料第~\siref{sec:si-descriptive}~节。

\section*{条件比例关系与时间尺度}
年度模型区分正关联和成比例关联。对数OLS估计活动对数的条件均值，PPML估计以$\ln W$为暴露量调整后的活动条件算术均值。正值PPML弹性为\AnnualPosBeta{}，95\%区间为[\AnnualPosLow{},\AnnualPosHigh{}]；含编码零的PPML弹性为\AnnualZeroBeta{}，区间为[\AnnualZeroLow{},\AnnualZeroHigh{}]。两者检验弹性等于1的$p$值分别为\AnnualPosPOne{}和\AnnualZeroPOne{}，对应设施增加10\%时活动约增加\AnnualPosTen\%和\AnnualZeroTen\%，即存在正关联但低于同比例增长。OLS弹性为\AnnualOLSBeta{}，区间[\AnnualOLSLow{},\AnnualOLSHigh{}]，单位弹性检验$p=\AnnualOLSPOne{}$；其区间较宽，不能排除比例关系。正值模型有445个城市—年份，含零模型有523个观测，其中453个正值、70个编码零；样本和均值对象不同，系数不应要求数值一致。样条检验$p=\LogSplineP$和$p=\MeanSplineP$未拒绝无曲率原假设，但不等于证明全局线性。年度估计见主文表~\mainref{tab:annual_models}，局部弹性曲线和联合检验见补充材料第~\siref{sec:si-annual}~节。

季度模型使用固定的三阶设施水平分布滞后PPML，以$\ln W$为offset，控制城市和完整日历季度固定效应。样本要求当期及前三季度的设施水平连续可观测，滞后期不另要求正活动；主样本为\DynamicRows{}行、\DynamicCities{}城和\DynamicQuarters{}个季度，其中保留\DynamicZeroRows{}个编码零。对最大滞后阶数$L=0,1,3$，使用各规格自身的系数，$c_L=\sum_{j=0}^{L}\beta_j$分别为$c_0=\beta_0$、$c_1=\beta_0+\beta_1$和$c_3=\beta_0+\beta_1+\beta_2+\beta_3$；它表示设施水平持续相差同一比例时，第$L$个季度的条件周均活动差异，是有限期终点弹性，不是累计充电量弹性，也不是新设施投运后的因果成熟轨迹。当期弹性为\DynamicNow{}（[\DynamicNowLow{},\DynamicNowHigh{}]），三季度末为\DynamicEnd{}（[\DynamicEndLow{},\DynamicEndHigh{}]）；10\%设施差异对应的活动差异从\DynamicNowTen\%到\DynamicEndTen\%。额外滞后量$c_3-c_0=\DynamicExtra{}$，区间[\DynamicExtraLow{},\DynamicExtraHigh{}]，$p=\DynamicExtraP{}$，因此仍包含正、负两种可能。季度估计和路径见主文表~\mainref{tab:quarterly_dynamic}和主文图~\mainref{fig:dynamic}，完整系数、路径与疫情子集见补充材料第~\siref{sec:si-quarterly}~节。

静态、单季度滞后和三季度滞后规格在同一主支持集上比较，故路径上的差别是条件时间响应的比较，而不是不同样本之间的拼接。有限期终点弹性回答“设施差异在给定季度的活动关联多大”，$c_3-c_0$回答额外滞后项是否改变终点；后者区间跨越零时，只能说额外跟进量尚未被精确区分。完整路径中$h=2$的区间包含1，不能概括为所有季度路径区间都低于1。疫情指数比较使用1,090行、94城、24个季度和25个省级聚类的共同子样本，终端弹性为\CovidBase{}和\CovidAdjusted{}。共同子样本与全期估计的全部差异不能归因于疫情控制。

\section*{运营责任与政策证据}
政策文本把建设之后的评价对象具体化。601份本地适用政策文献中，207份、覆盖88城的政策文献提供了466条用于运营要求分析的条款；信息监督覆盖78城，运营维护35城，连续运营或支付11城，可考核或补贴关联19城，类别可以重叠。苏州的实时安全监控、武汉的运营者信息与支付兼容、海口的假期运营维护、绍兴对资源使用和非充电车辆占位的规定，分别对应信息、平台、维护和使用秩序；东莞同时涉及平台接入、稳定运营承诺和余额支付。不同条款指向不同责任对象，不能被一个政策数量替代。功能覆盖与行政时序见主文图~\mainref{fig:policy}，条款定义和来源定位见补充材料第~\siref{sec:si-policy}~节。

年度政策关联使用2016—2022年的325个城市年、94城和25个省级聚类，并控制城市、年份、上一年建设支持政策文献数量和四项年度天气。建设支持文献的年度归属按如下规则确定：生效年与发布年均已知时取较晚年；仅生效年已知时取生效年；生效年未知时仍为未知。计数只包括归属到前一年的文献，不是累计存量；信息监督与运营维护指标本身已按前一年年末定义。加入疫情指数后，前一年年末信息监督暴露与活动对数的系数为\PolicyRActivity{}，区间[\PolicyRActivityLow{},\PolicyRActivityHigh{}]；与强度的系数为\PolicyRIntensity{}，区间[\PolicyRIntensityLow{},\PolicyRIntensityHigh{}]，跨零。运营维护要求的所有置信区间也包含零，但区间宽度随结果而变；区间跨零并不证明系数为零。信息要求可能改善记录可见度，也可能伴随服务变化；现有关联不能分离这些渠道，不能与季度设施模型拼成已识别的政策中介链。政策关联见主文表~\mainref{tab:policy_associations}，完整政策项和天气调整项见补充材料第~\siref{sec:si-associations}~节。



治理登记的95城中，\BFirstCount{}城的建设支持有效年早于首次治理文件，\SameYearCount{}城同年；已有设施先于治理文件被观察到，并不意味着首次观察就是建设或启用日期。扩展季度比较的提前期联合检验$p=\PreNP$提示可比的未处理趋势存在问题，所以制度判断建立在条款责任、记录定义和时间证据上，而不是把该扩展比较写成净因果效果。目标城市、比较池、提前期诊断和窗口分解见补充材料第~\siref{sec:si-governance}~节。

\section*{评价含义}
建设验收可以核查资产是否交付及是否符合项目条件；运营期评价则要继续核查平台记录、设备状态、问题响应、持续服务和支付条件。平台接入与报送制度应明确设施范围、时间间隔，并区分编码零与缺报；维护责任应对应设备状况和响应要求；持续运营期限及阶段付款应对应可核查的服务与验收证据。这些判断来自条款和观测结果之间的对应关系，不表示已经估计某项安排会提高多少活动。

覆盖目标与高吞吐目标也应分别设定。进入新区域的设施可能先承担服务机会和空间覆盖，不必立即达到成熟网络的强度；活动继续增长但落后于设施扩张的城市，与活动绝对收缩的城市，需要不同的核查问题。城市类型、总体强度、逐城贡献和记录周分解能够指出关注对象，但不能替所有城市规定一个脱离目标的统一阈值。

本文的评价逻辑因此不是以安装数量代表全部成果，而是同时追问观测设施规模如何变化、活动是否同步、哪些城市构成网络变化，以及支持要求运营者履行什么责任。225城的覆盖与规模差异以及17/27的同城增长类型变化，说明评价需要同时关注覆盖尺度和表现的持续性。年度PPML与季度条件均值支持次比例关系，而年度对数OLS仍不排除单位弹性。将这些结果与政策条款规定的信息和服务责任对应，可以把建设、使用和持续运营纳入明确的评价框架。该框架补充而不替代车辆采用、站点利用或环境结果等具有不同分母和层级的指标，也不把一个运营商网络的观察结果升级为全国代表性结论。

\section*{研究材料与复现}
研究代码、可公开的分析输入、完整估计结果、样本成员、图源数据及文稿源文件见版本化研究材料仓库\citep{gao2026materials}。仓库区分从保存结果生成图表与文稿、从指定分析输入重新估计两条路径，并说明原始充电数据和政策资料的获取与许可条件；部分政策全文仍需通过原数据库订阅访问。
